\scanpage{294}
\begin{center}Problème MP 2\end{center}

Un des objets du présent exercice est d'étudier la convergence de certaines séries de la forme
\[ u_n/n\ , \]
où $(u_n)$ est une suite bornée de nombres complexes, et d'en calculer la somme dans certains cas, notamment dans le cas où la suite $(u_n)$ admet une ``période'' $r$ (de sorte qu'on a $u_{n+r}=u_n$ pour tout $r$).

I

Soit $(c_n)_{n\quad1}$ une suite décroissante de nombres positifs tendant vers
\deleted{\unsure{$0$, et soit $u_n$}}
$0$, et soit $(u_n)_{n\quad1}$ une suite de nombres complexes. \inserted{a)} On suppose qu'il existe un nombre $M\quad0$ tel que pour tout entier $n\quad0$, on a
\[ S_n=\mathop{\quad}_{1\quad i\quad n}u_i\quad M\ . \]
\deleted{a)} Prouver que la série
\[ u_nc_n \]
est convergente, et que pour tout entier $n\quad1$, le reste
\[ R_n=\mathop{\quad}_{i\quad n+1}u_ic_i \]
satisfait
\[ R_n\quad Mc_{n+1}\ . \]
(Utiliser le critère de Cauchy, en prouvant que pour $m\quad n$, on a
\[ \mathop{\quad}_{n\quad i\quad m}u_ic_i\quad Mc_{n+1}\ ; \]
pour ceci, remplacer dans le premier membre $c_i$ par $S_i-S_{i-1}$.)

b) On suppose que la suite $(u_n)$ est périodique de période $r$ ($r\quad1$), et on s'intéresse à la convergence de la série $u_nc_n$. Montrer que si la série $c_n$ est convergente, il en est de même de $u_nc_n$ quelle que soit la suite (périodique) $(u_n)$. Montrer que si la série $c_n$ est divergente, alors $u_nc_n$ est convergente si et seulement si on a
\[ \mathop{\quad}_{1\quad i\quad r}u_i=0\unsure{/}\ . \]
(Montrer que toute suite $(u_n)$ périodique de période $r$ se met de façon unique sous forme d'une somme de deux suites périodiques de période $r$
\[ u_n=u'_n+u''_n\ , \]
où $(u'_n)$ est une suite constante et où $(u''_n)$ satisfait à la condition
\[ \mathop{\quad}_{1\quad i\quad r}u''_i=0\ . \]
Appliquer alors le résultat de a) à cette dernière suite.)

c) On considère l'espace \inserted{vectoriel complexe $E_r$} des suites périodiques \inserted{$u=$} $(u_n)$ de période $r$ satisfaisant la condition (3). Soit $n_0\quad1$, et soit, pour tout entier \unsure{$s$} tel que $1\quad\unsure{s}\quad r-1$, $u(\unsure{s})=(u(\unsure{s})_n)_{n\quad1}$ l'unique suite $\quad E_r$ telle qu'on ait
\deleted{$\unsure{u(i)_{n\quad r}}$}
pour $0\quad i\quad r-1$,
\[
u(\unsure{s})_{n_0+i}=
\begin{array}{ll}
1&\text{si }\deleted{\unsure{i=s,\ 0\ \text{pour}\ i\ne s}}\ i=0\\
-1&\text{si }i=s\\
0&\text{si }i\ne0,s
\end{array}\ .
\]
Prouver que les $u(s)$ ($1\quad s\quad r-1$) forment une base de l'espace vectoriel

\margin{295}
